% =====================================================================
%  Chapter 5 — Principle: Independence (Share Nothing)
%  Map: research/chapter-map.md (Chapter 5)
%  Shape: idea → origin → cost → what it opens up. Ideas, not how-to.
%  Rule: the author is never quoted in body text; short quotes only in
%  footnotes. Full statements live in interviews/.
% =====================================================================
\chaptersubtitle{Share Nothing}
\chapter{Independence}
\label{ch:actors}

A Silica program is not a tree of procedures calling procedures. It is a
society of independent actors, each with its own private state, that cooperate
by sending one another messages. No actor can reach into another's memory.
Independence is the normal case; sharing is a rare, deliberate exception.

\section{A borrowed idea, made native}

The idea comes from Erlang, and Silica makes no secret of the debt. Erlang's
designers built telephone switches that had to keep running for years, and they
found that independence was the foundation of reliability: if nothing is
shared, one failure cannot spread.

What Silica changes is where the idea can live. Erlang's actors depend on a
virtual machine and a garbage collector. Silica's are compiled directly to the
machine, with neither, so the same model can run on a server or on a
microcontroller small enough to fit on a badge.

\section{What SwErl taught}

Barney had tried once before to carry Erlang's model somewhere new. SwErl
brought it into Swift, and Barney counts it a good first exploration that taught
four lessons quickly.\footnote{Interview, 2026-11-13.}

The first was that supervisors were nearly impossible to build as a library.
There was no clean way for a crashing actor to report to its supervisor before
it disappeared, and without that report a supervisor cannot do a good job of
restarting what failed. The second was how little an actor really needs: a
stored state and a function, nothing more. There is no need for a hand-written
lock to make it wait for a message, and no need for some intermediate state in
which it has been spawned but is merely alive, waiting to be moved off a thread.
The third followed from the second: actors can be fairly simple to implement.
The fourth was a warning. Barney built something like traits, modelled on the
templates Erlang uses for its standard behaviours, but it was cumbersome, and
the results were uncomfortable.

A language can settle what a library cannot. Because Barney was writing
Silica's compiler, failure reporting could be built on whatever the chips and
operating systems themselves are able to report, and supervision became
possible from the start (Chapter~\ref{ch:failsoft}). Actors became a stored
state and a function, and traits found a cleaner form in the language itself.

\section{Any function can stand alone}

In Silica, being an actor is not a special kind of code. An ordinary function
that knows how to handle one message becomes an independent actor when the
program starts it as one. The runtime does the repetitive work of waiting for
the next message. The architecture of a Silica application is therefore its
society of actors, not its call graph, and its starting point is small: it
starts the society and steps aside.

\section{Don't wait by the phone}

When one actor needs something from another, Silica prefers that it send a
request and carry on, letting the answer arrive as a message of its own. The
documents compare this to texting: few people send a message and then freeze
until the reply comes.\footnote{\src{silica/docs/learn-silica.md}, line 913.}
Waiting is possible; it is simply not the first choice.

Messages are copies. Nothing a sender puts in a message remains shared with it,
so the sender's memory can be released at once and neither side can disturb the
other.

\section{Independence at scale}

Independence is only useful if it is cheap. When Silica's hosted runtime gave
every actor its own operating-system thread, the model could not reach the
numbers Erlang programmers take for granted. In September 2026 one of the AI
agents working on the project proposed scaling the promise back and offering
massive concurrency as a library instead. Barney refused. Actors bound one to
one to threads would be nothing but a fancy facade for a thread, and a language
that called them actors would be telling a lie.\footnote{``They would be false
and untruthful.'' Interview, 2026-11-13.} In October 2026 the runtime was
rebuilt so that many actors share each thread. A single machine then held a
million idle actors.

\section{The end of one size fits all}

One difference from Erlang is deliberate. Erlang's runtime moves work between
processor cores as it sees fit. Silica's never moves an actor on its own; the
program decides where its actors run, and the decision stays visible. The reason
is the principle of Chapter~\ref{ch:hidden}: nothing hidden. Any programmer can
write the balancing scheme that best fits the application at hand and the
larger problem it solves.

Behind that choice is a larger idea. One-size-fits-all solutions served
computing well for decades, because code was expensive to write and every
special case cost human effort. With code now generated quickly, Barney is no
longer convinced. The same reasoning reaches even the language: Silica and its
compiler can reasonably be forked when a need demands it, for instance to give
an application its own way of packaging and transporting data.

The obvious objection is that forks fragment a community, and that shared
libraries, knowledge and tools are what make a language ecosystem valuable.
Barney's answer is that the objection was right in the past. Communities still
form around shared knowledge, but the knowledge itself is changing. It is
becoming knowledge of how to shape the language to make a particular piece of
silicon do what is needed this time, rather than knowledge of how to bend a
library to fit.\footnote{``how do I get this piece of rock to do what I need it
to this time?'' Interview, 2026-11-13.} Libraries will still exist and still be
shared. The difference is that anyone can adapt one to the exact, direct need of
the problem in front of them, because of the language models now available.

\section{What it costs}

Thinking in independent actors is a different discipline from thinking in
procedures, and it must be learned. Copying messages costs time and memory that
sharing would save. And where Erlang offers automatic conveniences, such as
balancing load across cores, Silica asks the program to decide, which is work
someone must do.

\section{What it opens up}

Nearly all systems built today are static, and Barney sees that as an accident
of history: machines were slow and weak, and humans writing code were the
bottleneck. Cheap, independent actors invite a different kind of system. Barney
pictures millions of distributed nodes, each spawning actors, using them once
and discarding them, around a smaller number of actors that persist.

Erlang's networks of machines, a work of genius and the inspiration for all of
this, were designed for stable, safe environments. That is not the world
devices live in now. Devices are not always on, and networks are not always
there. A phone moves from Ethernet to Wi-Fi to Bluetooth in the middle of a
conversation with a single other device. Phones reboot; a phone call stops the
apps running on it. BEAM\_SEMP took Erlang's ideas a step further, toward
systems designed for exactly that churn, and Silica carries the idea on
(Chapter~\ref{ch:transport}).

\begin{itemize}
  \item \textbf{Architecture that matches intuition.} People already know how
        to think about friends who keep their own counsel and pass messages.
        A program shaped that way is easier to explain, to teach and to
        reason about.
  \item \textbf{One actor per thing.} When actors cost almost nothing, a
        design can give one to every connection, every sensor, every user or
        every task, and let each fail or succeed alone.
  \item \textbf{Every core, used safely.} Independent actors spread across
        processor cores without locks, without races and without the classic
        bugs of shared memory.
  \item \textbf{A model from the badge to the data centre.} The same design
        runs on a microcontroller and a server, so one way of thinking
        covers both.
  \item \textbf{Boundaries for security.} An actor that cannot touch another's
        memory is also an actor that cannot steal from it.
  \item \textbf{Systems built for churn.} Applications can be made of actors
        that come and go as devices sleep, wake, roam between networks and
        restart, instead of assuming a stable world.
\end{itemize}

Where this leads, in Barney's imagination, is beyond the operating system
itself. Chapter~\ref{ch:transport} ends there.
